% TOP_PROBLEM: {"id": 20001118, "title": "Square-difference sets, Paley cliques, and an exact external-defect variance identity", "category_id": 2, "category": {"id": 2, "name": "combinatorics", "display_name": "Combinatorics", "description": "Counting problems, graph theory, discrete structures.", "slug": "combinatorics", "order_index": 2, "created_at": "2026-07-31T15:26:25.671Z"}, "status": "partially_solved", "problem_number": "AIM-COMBINATORICS-0243", "set_id": 14}
% FIRST_POSTED: 2026-10-01
% ENTRY_KIND: statement_only
% UnsolvedMath ID 20001118; source catalogue code AIM-COMBINATORICS-0243
% !TeX program = lualatex
% Definition and sources only; this file is not an LLM solution attempt.
\documentclass[11pt,a4paper]{article}
\usepackage[margin=25mm]{geometry}
\usepackage{fontspec}
\setmainfont{lmroman10-regular.otf}[BoldFont=lmroman10-bold.otf,
 ItalicFont=lmroman10-italic.otf,BoldItalicFont=lmroman10-bolditalic.otf]
\usepackage{amsmath,amssymb,mathtools}
\usepackage{unicode-math}
\setmathfont{latinmodern-math.otf}
\AtBeginDocument{\let\setminus\smallsetminus}
\usepackage{xurl,hyperref}
\hypersetup{colorlinks=true,urlcolor=blue}
\setcounter{secnumdepth}{0}
\setlength{\parindent}{0pt}
\setlength{\parskip}{5pt}
\setlength{\emergencystretch}{3em}
\begin{document}
\section*{Square-difference sets, Paley cliques, and an exact external-defect variance identity}\label{20001118}
{\small UnsolvedMath ID 20001118 · AIM-COMBINATORICS-0243 · Combinatorics · AIM Workshop Problem Lists}
\subsection{Definitions and mathematical statement}
For an odd prime $p$, let
$Q_p=\{x^2:x\in\mathbb F_p^*\}$ be the set of nonzero quadratic
residues modulo $p$. Define
\[
 M(p)=\max\{|A|:A\subseteq\mathbb F_p,
                              A-A\subseteq Q_p\cup\{0\}\},
 \qquad A-A=\{a-b:a,b\in A\}.
\]
Thus the difference of every two distinct members, in each order,
must be a nonzero square. The diagonal differences are zero and
are explicitly exempted. This formulation removes any ambiguity
about whether zero is called a quadratic residue.

The substantive case is $p\equiv1\pmod4$. Then $-1$ is a square,
and the Paley graph on $\mathbb F_p$ has an undirected edge between
$a$ and $b$ exactly when $a-b\in Q_p$. The quantity $M(p)$ is
its clique number: the largest number of pairwise adjacent vertices.
The problem asks how large this number can be, ideally determining
its sharp dependence on $p$ or sharp uniform asymptotic bounds.
No assumption that $A$ is a subgroup, an interval, or a set of
residues of some fixed polynomial is imposed.

For completeness, when $p\equiv3\pmod4$, both $a-b$ and $b-a$
cannot be nonzero squares, so the same definition gives $M(p)=1$.
This explains why the nontrivial extremal question concerns primes
congruent to one modulo four. The problem involves pairwise additive
differences, even though eligibility of a difference is multiplicative.
\subsection{Short English statement}
How large can a subset of a prime field be when every nonzero difference between its elements is a square?
\subsection{Sources}
\begin{itemize}
\item[S1] ULAM.AI and the UnsolvedMath Contributors, supplied problems.json, record 20001118 (AIM-COMBINATORICS-0243), AIM Workshop Problem Lists. Catalogue identity and source transcription; CC BY 4.0.
\url{https://huggingface.co/datasets/ulamai/UnsolvedMath}.
\item[S2] American Institute of Mathematics, Recent trends in additive combinatorics, item 3.6. Original workshop question underlying AIM-COMBINATORICS-0243.
\url{https://aimath.org/WWN/additivecomb/additivecomb.pdf}.
\end{itemize}
\end{document}
